\documentclass[a4paper]{article}
\usepackage[fontset=fandol]{ctex}
\usepackage{amsmath, amssymb}
\usepackage{graphicx}
\usepackage[margin=2.5cm]{geometry}
\usepackage[most]{tcolorbox}
\usepackage{etoolbox}
\usepackage{listings}
\usepackage{hyperref}
\usepackage{booktabs}
\usepackage{subcaption}
\usepackage{float}
\usepackage{tikz}
\IfFileExists{pgfplots.sty}{\usepackage{pgfplots}\pgfplotsset{compat=1.18}}{}
\newtcolorbox{knowledgebox}[1]{enhanced, colback=blue!5!white, colframe=blue!75!black, colbacktitle=blue!75!black, coltitle=white, fonttitle=\bfseries, title=#1, attach boxed title to top left={yshift=-2mm, xshift=2mm}, boxrule=1pt, sharp corners}
\newtcolorbox{importantbox}[1]{enhanced, colback=yellow!10!white, colframe=yellow!80!black, colbacktitle=yellow!80!black, coltitle=black, fonttitle=\bfseries, title=#1, sharp corners}
\newtcolorbox{warningbox}[1]{enhanced, colback=red!5!white, colframe=red!75!black, colbacktitle=red!75!black, coltitle=white, fonttitle=\bfseries, title=#1, sharp corners}
\lstset{language=Python,basicstyle=\ttfamily\small,keywordstyle=\color{blue},stringstyle=\color{red!60!black},commentstyle=\color{green!60!black},breaklines=true,frame=single,numbers=left,numberstyle=\tiny\color{gray},captionpos=b,extendedchars=false}
\newcommand{\notetitle}{End-to-End ASR: CTC}
\newcommand{\noteauthors}{Codex 基于课程页、公开视频字幕、coverage ledger 与补充官方材料整理}
\newcommand{\notedate}{\today}
\newcommand{\videochannel}{WAVLab}
\newcommand{\videopublishdate}{2023-11-08}
\newcommand{\videoduration}{unknown}
\newcommand{\videourl}{https://www.youtube.com/watch?v=00GHpnRTAQE}
\newcommand{\videocoverpath}{cover.jpg}
\begin{document}
\begin{titlepage}
\centering
{\Large 课程笔记\par}
\vspace{1.2cm}
{\huge\bfseries \notetitle\par}
\vspace{0.8cm}
{\large \noteauthors\par}
\vspace{0.3cm}
{\large \notedate\par}
\vspace{1.2cm}
\ifdefempty{\videocoverpath}{{\small 公开视频封面不可用。\par}}{\includegraphics[width=0.82\textwidth,height=0.45\textheight,keepaspectratio]{\videocoverpath}\par}
\vfill
\begin{tcolorbox}[width=0.9\textwidth, colback=black!2!white, colframe=black!60, sharp corners]
\textbf{视频作者/频道}：\videochannel\par
\textbf{发布日期}：\videopublishdate\par
\textbf{视频时长}：\videoduration\par
\textbf{视频链接}：\href{\videourl}{\nolinkurl{\videourl}}
\end{tcolorbox}
\end{titlepage}
\tableofcontents
\newpage
\section{Modern encoder variants and Conformer bridge}
根据课程页主题、公开视频字幕以及本讲的 coverage ledger，Modern encoder variants and Conformer bridge 这一部分主要围绕 End-to-End ASR: CTC | End-to-End ASR: CTC；补完 attention 讨论后，引入 global/local / Branchformer-style encoder 设计，并把 Conformer 作为现代 ASR 编码器的代表。 展开。
写作上保持 coverage-first 原则：先把老师在这一阶段真正讲过的问题设定、模型部件、训练或推理逻辑讲清，再压缩成适合复习的结构。

\begin{figure}[H]
\centering
\includegraphics[width=0.92\textwidth]{figures/coverage_map.png}
\caption{本讲 coverage map：按 coverage\_units.jsonl 归纳出的核心主题与章节映射。}
\end{figure}

\begin{figure}[H]
\centering

\includegraphics[width=0.92\textwidth]{figures/timeline.png}
\caption{本讲 timeline：根据字幕时间范围归纳的主题推进顺序。}
\end{figure}

\subsection{End-to-End ASR: CTC}
从证据位置看，这一部分对应的课程证据区间是 2023-11-08。课程在这里强调的核心内容可以概括为：End-to-End ASR: CTC | End-to-End ASR: CTC。
CTC 的关键在于保持单调对齐假设，同时把对齐路径视作潜变量并在训练时进行求和。

在本章结构中，这一单元被并入“Modern encoder variants and Conformer bridge”，目的是把同一阶段的定义、机制和权衡集中讲清。
\[
P(Y \mid X) = \sum_{\pi \in \mathcal{B}^{-1}(Y)} \prod_{t=1}^{T} P(\pi_t \mid X)
\]

符号说明如下：
\begin{itemize}
\item $X$：输入语音特征序列；
\item $Y$：输出标签序列；
\item $\pi$：带有 blank 的对齐路径；
\item $\mathcal{B}$：把路径压缩为最终标签序列的映射。
\end{itemize}

\subsection{补完 attention 讨论后，引入 global/local / Branchformer-style encoder 设计，并把 Conformer 作为现代 ASR 编码器的代表}
从证据位置看，这一部分对应的课程证据区间是 00:00:01,040--00:07:24,120。课程在这里强调的核心内容可以概括为：补完 attention 讨论后，引入 global/local / Branchformer-style encoder 设计，并把 Conformer 作为现代 ASR 编码器的代表。。
与单调模型不同，attention encoder-decoder 更强调全局条件化能力，但也因此引入了对齐自由度和部署层面的权衡。

记录系统中的附加说明指出：CS224S supplement anchor: Conformer / Branchformer / Whisper-style encoders; useful if you want a modern encoder comparison.

在本章结构中，这一单元被并入“Modern encoder variants and Conformer bridge”，目的是把同一阶段的定义、机制和权衡集中讲清。
\begin{knowledgebox}{CS224S 补充}
CS224S supplement anchor: Conformer / Branchformer / Whisper-style encoders; useful if you want a modern encoder comparison.
\end{knowledgebox}

\subsection{实现视角与伪代码}
为了把这一讲的算法流程落到工程实现层面，下面给出一个教材化的伪代码（pseudocode）摘要。它不是逐行复刻讲者的实现，而是把 lecture 中反复强调的计算顺序、状态更新与解码决策压缩成可复查的程序骨架。
\begin{lstlisting}[caption={CTC training sketch}]
for utterance in dataset:
    logits = encoder(utterance)
    alpha = forward_dp(logits, target_with_blank)
    beta = backward_dp(logits, target_with_blank)
    loss = -log(sum_valid_paths(alpha, beta))
    update_model(loss)
\end{lstlisting}
把这一段放在正文里有两个目的：第一，让读者把抽象的模型或动态规划递推与真实实现步骤对齐；第二，为后续阅读 ESPnet、Kaldi 或其它 toolkit 时建立可迁移的 mental model。

\subsection{本章小结}
Modern encoder variants and Conformer bridge 这一部分把 End-to-End ASR: CTC 中相邻的教学单元串成了一条连续叙事：既交代了这一阶段要解决的问题，也交代了方法、推理或工程权衡为什么要这样组织。

\section{CTC formulation}
根据课程页主题、公开视频字幕以及本讲的 coverage ledger，CTC formulation 这一部分主要围绕 CTC 的基本 formulation、标签序列与输入帧序列之间的长度不匹配问题。 展开。
写作上保持 coverage-first 原则：先把老师在这一阶段真正讲过的问题设定、模型部件、训练或推理逻辑讲清，再压缩成适合复习的结构。

\subsection{CTC 的基本 formulation、标签序列与输入帧序列之间的长度不匹配问题}
从证据位置看，这一部分对应的课程证据区间是 00:07:21,000--00:14:10,000。课程在这里强调的核心内容可以概括为：CTC 的基本 formulation、标签序列与输入帧序列之间的长度不匹配问题。。
CTC 的关键在于保持单调对齐假设，同时把对齐路径视作潜变量并在训练时进行求和。

在本章结构中，这一单元被并入“CTC formulation”，目的是把同一阶段的定义、机制和权衡集中讲清。
\[
P(Y \mid X) = \sum_{\pi \in \mathcal{B}^{-1}(Y)} \prod_{t=1}^{T} P(\pi_t \mid X)
\]

符号说明如下：
\begin{itemize}
\item $X$：输入语音特征序列；
\item $Y$：输出标签序列；
\item $\pi$：带有 blank 的对齐路径；
\item $\mathcal{B}$：把路径压缩为最终标签序列的映射。
\end{itemize}

\subsection{本章小结}
CTC formulation 这一部分把 End-to-End ASR: CTC 中相邻的教学单元串成了一条连续叙事：既交代了这一阶段要解决的问题，也交代了方法、推理或工程权衡为什么要这样组织。

\section{Blank symbol and path marginalization}
根据课程页主题、公开视频字幕以及本讲的 coverage ledger，Blank symbol and path marginalization 这一部分主要围绕 blank symbol、路径求和与 many-to-one alignments 如何把 CTC 变成可优化目标。 展开。
写作上保持 coverage-first 原则：先把老师在这一阶段真正讲过的问题设定、模型部件、训练或推理逻辑讲清，再压缩成适合复习的结构。

\subsection{blank symbol、路径求和与 many-to-one alignments 如何把 CTC 变成可优化目标}
从证据位置看，这一部分对应的课程证据区间是 00:14:06,079--00:20:51,720。课程在这里强调的核心内容可以概括为：blank symbol、路径求和与 many-to-one alignments 如何把 CTC 变成可优化目标。。
CTC 的关键在于保持单调对齐假设，同时把对齐路径视作潜变量并在训练时进行求和。

在本章结构中，这一单元被并入“Blank symbol and path marginalization”，目的是把同一阶段的定义、机制和权衡集中讲清。
\[
P(Y \mid X) = \sum_{\pi \in \mathcal{B}^{-1}(Y)} \prod_{t=1}^{T} P(\pi_t \mid X)
\]

符号说明如下：
\begin{itemize}
\item $X$：输入语音特征序列；
\item $Y$：输出标签序列；
\item $\pi$：带有 blank 的对齐路径；
\item $\mathcal{B}$：把路径压缩为最终标签序列的映射。
\end{itemize}

\subsection{本章小结}
Blank symbol and path marginalization 这一部分把 End-to-End ASR: CTC 中相邻的教学单元串成了一条连续叙事：既交代了这一阶段要解决的问题，也交代了方法、推理或工程权衡为什么要这样组织。

\section{Dynamic programming for CTC}
根据课程页主题、公开视频字幕以及本讲的 coverage ledger，Dynamic programming for CTC 这一部分主要围绕 forward-backward / Viterbi 这类 dynamic programming 如何服务于 CTC 训练。 展开。
写作上保持 coverage-first 原则：先把老师在这一阶段真正讲过的问题设定、模型部件、训练或推理逻辑讲清，再压缩成适合复习的结构。

\subsection{forward-backward / Viterbi 这类 dynamic programming 如何服务于 CTC 训练}
从证据位置看，这一部分对应的课程证据区间是 00:20:48,320--00:28:11,080。课程在这里强调的核心内容可以概括为：forward-backward / Viterbi 这类 dynamic programming 如何服务于 CTC 训练。。
CTC 的关键在于保持单调对齐假设，同时把对齐路径视作潜变量并在训练时进行求和。

在本章结构中，这一单元被并入“Dynamic programming for CTC”，目的是把同一阶段的定义、机制和权衡集中讲清。
\[
P(Y \mid X) = \sum_{\pi \in \mathcal{B}^{-1}(Y)} \prod_{t=1}^{T} P(\pi_t \mid X)
\]

符号说明如下：
\begin{itemize}
\item $X$：输入语音特征序列；
\item $Y$：输出标签序列；
\item $\pi$：带有 blank 的对齐路径；
\item $\mathcal{B}$：把路径压缩为最终标签序列的映射。
\end{itemize}

\subsection{本章小结}
Dynamic programming for CTC 这一部分把 End-to-End ASR: CTC 中相邻的教学单元串成了一条连续叙事：既交代了这一阶段要解决的问题，也交代了方法、推理或工程权衡为什么要这样组织。

\section{CTC decoding and LM fusion}
根据课程页主题、公开视频字幕以及本讲的 coverage ledger，CTC decoding and LM fusion 这一部分主要围绕 CTC 的 decoding、beam search 与语言模型融合，及其在实际 ASR 系统中的角色。 展开。
写作上保持 coverage-first 原则：先把老师在这一阶段真正讲过的问题设定、模型部件、训练或推理逻辑讲清，再压缩成适合复习的结构。

\subsection{CTC 的 decoding、beam search 与语言模型融合，及其在实际 ASR 系统中的角色}
从证据位置看，这一部分对应的课程证据区间是 00:28:09,760--00:36:35,599。课程在这里强调的核心内容可以概括为：CTC 的 decoding、beam search 与语言模型融合，及其在实际 ASR 系统中的角色。。
CTC 的关键在于保持单调对齐假设，同时把对齐路径视作潜变量并在训练时进行求和。

记录系统中的附加说明指出：CS224S supplement anchor: compare with the lecture on Conformer / fine-tuning base speech models.

在本章结构中，这一单元被并入“CTC decoding and LM fusion”，目的是把同一阶段的定义、机制和权衡集中讲清。
\[
P(Y \mid X) = \sum_{\pi \in \mathcal{B}^{-1}(Y)} \prod_{t=1}^{T} P(\pi_t \mid X)
\]

符号说明如下：
\begin{itemize}
\item $X$：输入语音特征序列；
\item $Y$：输出标签序列；
\item $\pi$：带有 blank 的对齐路径；
\item $\mathcal{B}$：把路径压缩为最终标签序列的映射。
\end{itemize}
\begin{knowledgebox}{CS224S 补充}
CS224S supplement anchor: compare with the lecture on Conformer / fine-tuning base speech models.
\end{knowledgebox}

\subsection{本章小结}
CTC decoding and LM fusion 这一部分把 End-to-End ASR: CTC 中相邻的教学单元串成了一条连续叙事：既交代了这一阶段要解决的问题，也交代了方法、推理或工程权衡为什么要这样组织。

\section{Data and robustness}
根据课程页主题、公开视频字幕以及本讲的 coverage ledger，Data and robustness 这一部分主要围绕 external data、误标记 / 数据噪声与 CTC 系统鲁棒性的关系。 展开。
写作上保持 coverage-first 原则：先把老师在这一阶段真正讲过的问题设定、模型部件、训练或推理逻辑讲清，再压缩成适合复习的结构。

\subsection{external data、误标记 / 数据噪声与 CTC 系统鲁棒性的关系}
从证据位置看，这一部分对应的课程证据区间是 00:36:31,800--00:43:39,040。课程在这里强调的核心内容可以概括为：external data、误标记 / 数据噪声与 CTC 系统鲁棒性的关系。。
CTC 的关键在于保持单调对齐假设，同时把对齐路径视作潜变量并在训练时进行求和。

在本章结构中，这一单元被并入“Data and robustness”，目的是把同一阶段的定义、机制和权衡集中讲清。
\[
P(Y \mid X) = \sum_{\pi \in \mathcal{B}^{-1}(Y)} \prod_{t=1}^{T} P(\pi_t \mid X)
\]

符号说明如下：
\begin{itemize}
\item $X$：输入语音特征序列；
\item $Y$：输出标签序列；
\item $\pi$：带有 blank 的对齐路径；
\item $\mathcal{B}$：把路径压缩为最终标签序列的映射。
\end{itemize}

\subsection{本章小结}
Data and robustness 这一部分把 End-to-End ASR: CTC 中相邻的教学单元串成了一条连续叙事：既交代了这一阶段要解决的问题，也交代了方法、推理或工程权衡为什么要这样组织。

\section{Decode-time efficiency}
根据课程页主题、公开视频字幕以及本讲的 coverage ledger，Decode-time efficiency 这一部分主要围绕 更快的决策路径、beam search 约束以及对实时推理的影响。 展开。
写作上保持 coverage-first 原则：先把老师在这一阶段真正讲过的问题设定、模型部件、训练或推理逻辑讲清，再压缩成适合复习的结构。

\subsection{更快的决策路径、beam search 约束以及对实时推理的影响}
从证据位置看，这一部分对应的课程证据区间是 00:43:37,200--00:50:45,920。课程在这里强调的核心内容可以概括为：更快的决策路径、beam search 约束以及对实时推理的影响。。
Search 这一讲把前面所有模型真正接到了 decoding pipeline 上，也就是把训练得到的分数转成最终转写结果。

在本章结构中，这一单元被并入“Decode-time efficiency”，目的是把同一阶段的定义、机制和权衡集中讲清。

\subsection{本章小结}
Decode-time efficiency 这一部分把 End-to-End ASR: CTC 中相邻的教学单元串成了一条连续叙事：既交代了这一阶段要解决的问题，也交代了方法、推理或工程权衡为什么要这样组织。

\section{Alignment extensions and masked prediction}
根据课程页主题、公开视频字幕以及本讲的 coverage ledger，Alignment extensions and masked prediction 这一部分主要围绕 mask prediction、alignment issues 与 CTC 相关的扩展视角。 展开。
写作上保持 coverage-first 原则：先把老师在这一阶段真正讲过的问题设定、模型部件、训练或推理逻辑讲清，再压缩成适合复习的结构。

\subsection{mask prediction、alignment issues 与 CTC 相关的扩展视角}
从证据位置看，这一部分对应的课程证据区间是 00:50:43,280--00:57:34,240。课程在这里强调的核心内容可以概括为：mask prediction、alignment issues 与 CTC 相关的扩展视角。。
CTC 的关键在于保持单调对齐假设，同时把对齐路径视作潜变量并在训练时进行求和。

在本章结构中，这一单元被并入“Alignment extensions and masked prediction”，目的是把同一阶段的定义、机制和权衡集中讲清。
\[
P(Y \mid X) = \sum_{\pi \in \mathcal{B}^{-1}(Y)} \prod_{t=1}^{T} P(\pi_t \mid X)
\]

符号说明如下：
\begin{itemize}
\item $X$：输入语音特征序列；
\item $Y$：输出标签序列；
\item $\pi$：带有 blank 的对齐路径；
\item $\mathcal{B}$：把路径压缩为最终标签序列的映射。
\end{itemize}

\subsection{本章小结}
Alignment extensions and masked prediction 这一部分把 End-to-End ASR: CTC 中相邻的教学单元串成了一条连续叙事：既交代了这一阶段要解决的问题，也交代了方法、推理或工程权衡为什么要这样组织。

\section{Recap and assignment note}
根据课程页主题、公开视频字幕以及本讲的 coverage ledger，Recap and assignment note 这一部分主要围绕 CTC 整讲收束、优缺点与作业 / 实验提示。 展开。
写作上保持 coverage-first 原则：先把老师在这一阶段真正讲过的问题设定、模型部件、训练或推理逻辑讲清，再压缩成适合复习的结构。

\subsection{CTC 整讲收束、优缺点与作业 / 实验提示}
从证据位置看，这一部分对应的课程证据区间是 00:57:31,760--01:04:35,119。课程在这里强调的核心内容可以概括为：CTC 整讲收束、优缺点与作业 / 实验提示。。
CTC 的关键在于保持单调对齐假设，同时把对齐路径视作潜变量并在训练时进行求和。

在本章结构中，这一单元被并入“Recap and assignment note”，目的是把同一阶段的定义、机制和权衡集中讲清。
\[
P(Y \mid X) = \sum_{\pi \in \mathcal{B}^{-1}(Y)} \prod_{t=1}^{T} P(\pi_t \mid X)
\]

符号说明如下：
\begin{itemize}
\item $X$：输入语音特征序列；
\item $Y$：输出标签序列；
\item $\pi$：带有 blank 的对齐路径；
\item $\mathcal{B}$：把路径压缩为最终标签序列的映射。
\end{itemize}

\subsection{本章小结}
Recap and assignment note 这一部分把 End-to-End ASR: CTC 中相邻的教学单元串成了一条连续叙事：既交代了这一阶段要解决的问题，也交代了方法、推理或工程权衡为什么要这样组织。

\section{总结与延伸}
End-to-End ASR: CTC 这一讲在整门《Speech Recognition and Understanding》课程中的位置，是把 E2E ASR: CTC 讲成一个可复用的建模模块。
从教材化角度看，本讲最重要的不是记住单一术语，而是理解它如何嵌入完整的 automatic speech recognition pipeline：输入语音如何被表示、模型如何被训练、解码如何得到最终转写，以及这些设计分别带来什么假设和权衡。
由于公开源以课程页和公开视频字幕为主、缺少官方 slide PDF，本章在正文里尽量把术语、方法和工程语境解释完整，并把缺失项留在 omission 和 manifest 里，而不是凭空补写。
\end{document}
